\documentclass[
	11pt,
	aspectratio=169,
]{beamer}

\graphicspath{{Images/}{./}}

\usepackage{booktabs}
\usepackage[brazil]{babel}   
\usepackage[utf8]{inputenc}

\usepackage[autostyle]{csquotes}
\usepackage[style=authortitle-icomp,backend=biber]{biblatex}
\addbibresource{presentation.bib}

\usepackage[]{hyperref}
\hypersetup{colorlinks=true}

\usetheme{Palo Alto}
\usefonttheme{default} 
\usepackage{palatino} 
\usepackage[default]{opensans} 
\useinnertheme{circles}

\setbeamertemplate{navigation symbols}{}


\title[Engenharia de Plataforma]{Engenharia de Plataforma?}
\author[Murilo Alves]{Murilo Henrique Alves} 
\institute[UTFPR]{Universidade Tecnológica Federal do Paraná} 
\date[\today]{\today}

\begin{document}

\begin{frame}
	\titlepage
\end{frame}

\begin{frame}
	\frametitle{O que veremos?}
	
	\tableofcontents
	%\tableofcontents[pausesections]
\end{frame}


\section{Engenharia de Plataforma}

\begin{frame}
	\frametitle{O que é?}

	A Engenharia de plataforma é uma disciplina da engenharia de software que tem como objetivo
  projetar e desenvolver cadeias de ferramentas e fluxos de trabalho de autoatendimento,
  especialmente no contexto "cloud-native" \footcite{galante}.
\end{frame}

\begin{frame}
	\frametitle{De onde surgiu?}

	Engenharia de Plataforma surge à partir de necessidades específicas de times de engenharia relacionado com
	a aplicação na prática da cultura DevOps.

	\bigskip

	A cultura DevOps tem como uma de suas primícias a frase: "You build it, you run it", e aplicar isso na prática
	pode ser algo simples para organizações complexas como Google, AirBnb e Amazon, porém, é algo pouco realista pensar
	em aplicar a verdadeira cultura DevOps na prática na maioria dos times de engenharia nas empresas ''comuns''.

\end{frame}

\begin{frame}
  \frametitle{Ferramentas de DevOps}
	\framesubtitle{Um pequeno apanhado de tecnologias utilizadas para um workflow de entrega contínua em produção}

	\begin{itemize}
		\item Terraform (Infra as Code)
		\item Docker
		\item Helm Charts
		\item Kubernets
		\item GitOps
		\item CI (circleci, github actions, azure devops, bitbucket pipelines, jenkins, gitlab, tekton)
		\item CD (argocd, fluxcd, harnesscd, tekton)
		\item Secrets Manager (Hashicorp Vault, AWS SM)
		\item AWS (SQS, SNS, EC2, ECS, Lambda, CloudFront)
	\end{itemize}

\end{frame}

\begin{frame}

	Segundo Matthew Skelton \footcite{matthew}, quando empresas comuns tentam implementar
	na prática uma cultura DevOps, geralmente elas acabam gerando uma série de anti-padrões,
	devido à vários fatores como:
    \begin{itemize}
      \item é improvável que a maioria das empresas tenham acesso ao mesmo conjunto de talentos
      e ao mesmo nível de recursos que podem investir em melhoria do fluxo de trabalho
      e a experiência de desenvolvimento.
    \end{itemize}

\end{frame}

\begin{frame}
	\frametitle{O que faz?}

	As equipes de engenharia de plataforma geralmente fornecem o que é conhecido como
"Internal Developer Platform" (IDP), que visa atender a todas as necessidades
operacionais ao longo do ciclo de vida de um aplicativo.

\end{frame}

\section{Internal Developer Platform}

\begin{frame}
	\frametitle{O que é?}

	 Uma IDP engloba um conjunto de ferramentas e tecnologias, integradas de uma maneira à reduzir
	 a carga cognitiva nos desenvolvedores e ajudando a mantê-los em seus contextos originais de produto.
	\bigskip

	\begin{quote}
		Uma plataforma interna de desenvolvedor se concentra nas práticas de desenvolvimento interno de uma empresa.
	Você define um conjunto de caminhos de desenvolvimento recomendados e suportados para a produção e
	"pavimenta" incrementalmente um caminho através deles com uma plataforma interna.\\
		--- Kulla-Mader \& Lantz \footcite{kulla}
	\end{quote}

\end{frame}

\begin{frame}
	\frametitle{Golden Paths e Paved Roads}
  
  Esses caminhos de desenvolvimento recomendados são comunmente chamados de "Golden Paths" ou "Paved Roads".

  \bigskip

  Uma "Paved Road" consiste em um conjunto de processos e/ou ferramentas que são necessárias nos fluxos de
  entrega contínua mas que não são necessariamente de propriedade dos times de produtos, então os engenheiros
  de plataforma definem quais são esses processos e ferramentas e fornecem a cola para unir tudo em uma experiência
  de autoatendimento consistente.
\end{frame}

\begin{frame}
  \frametitle{A cola}

  Essa cola é a plataforma interna. Nas palavras de Evan Bottcher (Thoughworks) \footcite{evan}, plataformas são \begin{quote}
    uma base de APIs, ferramentas, serviços, conhecimento e suporte de autoatendimento, organizados como um produto interno atraente.
    Equipes de entrega autônomas podem usar a plataforma para entregar recursos de produtos em um ritmo mais acelerado,
    com coordenação reduzida.
  \end{quote}
  
\end{frame}

\begin{frame}[allowframebreaks]
	\frametitle{Referências}
	
	\begin{thebibliography}{99}
		\footnotesize
		
    	\printbibliography
    
	\end{thebibliography}
\end{frame}

\begin{frame}[plain]
	{\Huge The End}
\end{frame}

\end{document}